\documentclass[12pt]{article}
\usepackage[utf8]{inputenc}
\usepackage[T1]{fontenc}
\usepackage{hyperref}
\title{Yokuts Language Resources for K-12 Curriculum Development}
\author{Compiled by Claude Sonnet 5.5 (claude-sonnet-5-5) for Andr\'as Kornai}
\date{}
\begin{document}
\maketitle

\section{Introduction}
Yokuts is a family of related California languages and dialects (Yokutsan), not a single language. The ISO code yok covers it. Yawelmani, a Southern Valley Yokuts dialect, is spoken near the Tule River Reservation. A 2011 estimate by Victor Golla gave up to 25 fluent and semi-fluent speakers. Most active documentation work in the results concerns Chukchansi, a different Yokuts language.

\section{Primary Educational Resources}
No school curricula for Yawelmani were found. Two programme leads turned up:
\begin{itemize}
\item Fresno State's Chukchansi Yokuts Revitalization Project, founded in 2009. Its class materials, community grammar and bilingual dictionary are said to be still under revision: \url{https://cah.fresnostate.edu/linguistics/research/ilg.html}
\item A 1993 Master-Apprentice pilot that included Yawelmani was reported as effective for teaching conversational Yawelmani (see the Wikipedia page below).
\end{itemize}

\section{Linguistic and Reference Materials}
\begin{itemize}
\item Overview of the Yokuts languages: \url{https://cla.berkeley.edu/languages/yokuts.php}
\item Yawelmani Yokuts: \url{https://en.wikipedia.com/wiki/Yawelmani}
\item Stanley Newman, ``The Yawelmani Dialect of Yokuts'' (1932): \url{https://glottolog.org/resource/reference/id/556815}
\item Newman's 1946 sketch appeared in Hoijer, \emph{Linguistic Structures of Native America}.
\item Mary Haas's thumbnail sketch of Yawelmani is held in the Berkeley archives (Haas.063). No link was found.
\item Gamble, \emph{Wikchamni Grammar} (1978) and \emph{Yokuts Texts} (1994). These are printed books.
\item Collord, a Chukchansi grammar (1968 dissertation).
\item Article on Chukchansi in the Journal of the IPA: \url{https://www.cambridge.org/core/journals/journal-of-the-international-phonetic-association/article/abs/chukchansi-yokuts/EF895162E70C7C3B00350181BDBEBC69}
\end{itemize}
I found no current Yawelmani dictionary.

\section{Religious and Community Resources}
None found. The search suggested contacting the Tule River Tribe's language programme, but I found no page for it.

\section{Digital and Audio Resources}
None found for Yawelmani. Fresno State's work also includes Choinimni and Chalon dictionary compilation: \url{https://cah.fresnostate.edu/linguistics/research/ilg.html}

\section{Curriculum Development Context}
Any curriculum has to start by choosing a variety. Chukchansi has the most active revitalization work and Yawelmani the older grammatical documentation. The older academic grammars would need adapting before classroom use, and community consultation is essential.
\end{document}